% =====================================================================
% Chapter 1, part 1: Motivation, Problem Statement, Contributions,
% Structure. Transcribed from the reviewed REMAINING_SECTIONS.html draft.
% =====================================================================

\section{Motivation}
\label{sec:motivation}

Parkinson's disease is the second most common neurodegenerative disorder,
and its diagnosis remains clinical: a neurologist examines the patient's
movement against consensus criteria \cite{postuma2015}. Among the motor
signs, \textbf{gait} holds a special place. It deteriorates early, it
disables profoundly, and it is measurable: patients walk more slowly, with
shorter steps, reduced arm swing and, characteristically, greater
stride-to-stride variability and marked difficulty with turns and with
rising from a chair \cite{mirelman2019,hausdorff2009}. A technology that
measured gait objectively, repeatedly and outside the clinic could support
earlier referral, longitudinal tracking and therapy adjustment.

The instruments available today each concede something. Optical motion
capture is the laboratory gold standard but is expensive, confined to
instrumented rooms and requires markers. Wearable sensors travel with the
patient but must be worn, charged and accepted, a real burden for elderly
users. Cameras are cheap but raise immediate privacy objections in the
home. \textbf{Radar} concedes none of these: it is contactless, works
through clothing and in darkness, records no camera image, and its
micro-Doppler return integrates the motion of every moving part of the body
into one time-frequency image, the \textbf{micro-Doppler signature}, from
which the velocities of the torso and the limbs can be extracted
\cite{seifert2020,gurbuz2024}. In that image the torso traces a slow band
and each foot a fast periodic arch.

This thesis works on data from a clinically validated multi-node FMCW radar
network built precisely for gait monitoring, whose parameter extraction has
been validated in both healthy and Parkinson's populations
\cite{lopezdelgado2026}. The cohort contains 58 subjects, 33 healthy
controls and 25 Parkinson's patients, each walking a short corridor in a
protocol modelled on the clinical Timed-Up-and-Go test. The question this
thesis asks of that data is the natural next one: can the disease itself be
read from these signatures?

\section{Problem Statement}
\label{sec:problem}

Radar has already proven itself around gait, but on different questions.
One line of work \textbf{measures} gait: it validates which biomechanical
parameters, such as step time, cadence and stride length, can be extracted
reliably from the radar return \cite{seifert2020,lopezdelgado2026}. Another
line \textbf{classifies} walkers into coarse groups, almost always young
versus elderly adults, and reports accuracies above 90\,\%
\cite{hayashi2021,hoshiga2021}. Between the two sits the clinical question
this thesis takes up: whether Parkinson's disease itself can be recognised
in the radar signature of a walk, tested the only way a screening tool is
used in practice, on people the classifier has never seen. To the best of
our knowledge, no published work answers it.

How a classifier is tested decides what its score means, and this is where
the field is weakest. A micro-Doppler recording carries a strong personal
signature: published systems recognise individual \emph{people} from their
walk \cite{papanastasiou2020,ni2020}. If recordings of the same person
appear in both the training and the test data, a classifier can therefore
score highly by recognising the person and recalling their label, having
learned nothing about the disease. The field's own record shows the danger
in another form: a network trained on micro-Doppler images and reaching
97.8\,\% accuracy was shown
by its own authors to have keyed on background differences between
recording sites rather than on gait \cite{hayashi2021}. Other regularities
offer the same shortcut, because Parkinson's cohorts are typically older
than their controls, and patients take longer to complete a walking test.

The problem this thesis addresses is therefore twofold. First, to build and
evaluate a Parkinson's-versus-control classifier on radar micro-Doppler
signatures with every scored person excluded from training. Second, to do
so under an evaluation design in which every score is defensible: each
confound is measured, given an explicit baseline score that any real model
must beat, and controlled rather than ignored.

\section{Contributions}
\label{sec:contributions}

The thesis makes five contributions:

\begin{itemize}
  \item \textbf{The first subject-independent evaluation, to our knowledge,
        of Parkinson's-versus-control classification from radar
        micro-Doppler gait signatures.} The evaluation covers classical
        models and four deep model families, including AlexNet, the
        architecture behind the field's best-known result
        \cite{hayashi2021}.
  \item \textbf{A confound-aware evaluation protocol.} Every model is
        trained and scored with leave-one-subject-out validation on
        fixed-length windows, is judged against deliberately simple
        baseline predictors, such as one that knows only how long each
        recording lasts, and is reported with a subject-level confidence
        interval.
  \item \textbf{An honest measurement of how much signal the data holds.}
        Under this subject-independent evaluation, no classifier can be
        distinguished from the duration-only baseline. The finding is
        presented as a property of the data, with its uncertainty, not as a
        modelling failure.
  \item \textbf{An explanation of where the scores come from.} Attention
        analysis shows what the image networks actually read; an
        intervention experiment removes what they attend to; and a protocol
        experiment reproduces the literature's numbers on this very dataset
        and shows them collapsing on unseen subjects.
  \item \textbf{A clinically usable finding and a reusable protocol.} Radar
        reliably measures how long the walk and its phases take, in effect
        an automated Timed-Up-and-Go test, and the evaluation protocol can
        serve as the design document for the richer datasets the field
        needs next.
\end{itemize}

\section{Structure of this Document}
\label{sec:structure}

The rest of this document is organised as follows:

\begin{itemize}
  \item \textbf{Chapter~\ref{ch:soa}, State of the Art:}
    \begin{itemize}
      \item Radar micro-Doppler as a contactless gait sensing modality;
      \item The classification studies closest to this task;
      \item The validation pitfalls that recur throughout the field.
    \end{itemize}
  \item \textbf{Chapter~\ref{ch:methods}, Materials and Methods:}
    \begin{itemize}
      \item The dataset and the acquisition system;
      \item Preprocessing, windowing and the extracted features;
      \item The classical and deep model families;
      \item The leave-one-subject-out protocol, the confound controls and
            the evaluation metrics.
    \end{itemize}
  \item \textbf{Chapter~\ref{ch:results}, Results:}
    \begin{itemize}
      \item The exploratory characterisation that constrained the design;
      \item The baseline scores every model must beat;
      \item The classical and deep model campaigns and the input variants;
      \item Decision-level performance at a fixed threshold.
    \end{itemize}
  \item \textbf{Chapter~\ref{ch:discussion}, Discussion, Conclusions and Future Work:}
    \begin{itemize}
      \item The experiments that show where the gap to the literature
            comes from;
      \item What the models actually read, and the intervention that tests
            it;
      \item The positive findings, the limitations and the conclusions;
      \item The path to stronger results.
    \end{itemize}
  \item \textbf{Annexes:}
    \begin{itemize}
      \item Ethical, social, economic and environmental aspects;
      \item The project budget.
    \end{itemize}
\end{itemize}

\paragraph{Code and reproducibility.} The complete pipeline, the frozen
pre-registered configuration, the analysis scripts that produce every figure
and table, and the \LaTeX{} source of this document are publicly available at
\url{\repourl}. The clinical recordings themselves are not redistributed:
they remain with the originating team under the ethics approval described in
Annex~A, and the repository documents how to run the pipeline on them.
